\documentclass[ecta,nameyear,draft]{econsocart}
\usepackage{xr-hyper}
\usepackage{amsmath,amsfonts,amssymb,amsthm,graphicx,booktabs,array,longtable}
\usepackage{algorithm,algpseudocode,bm,enumitem}
\usepackage[T1]{fontenc}
\usepackage{mathpazo}
\usepackage{microtype}
\usepackage{needspace}
\RequirePackage[colorlinks,citecolor=blue,linkcolor=blue,urlcolor=blue]{hyperref}
\startlocaldefs
\newtheorem{theorem}{Theorem}
\newtheorem{proposition}[theorem]{Proposition}
\newtheorem{lemma}[theorem]{Lemma}
\newtheorem{corollary}[theorem]{Corollary}
\newtheorem{assumption}{Assumption}
\theoremstyle{definition}
\newtheorem{definition}{Definition}
\theoremstyle{remark}
\newtheorem{remark}{Remark}
\newtheorem{example}{Example}
\newcommand{\E}{\mathbb E}
\newcommand{\R}{\mathbb R}
\newcommand{\HH}{\mathcal H}
\newcommand{\LL}{\mathcal L}
\newcommand{\Act}{\mathcal A}
\newcommand{\sg}{\operatorname{sg}}
\newcommand{\tr}{\operatorname{tr}}
\newcommand{\Lip}{\operatorname{Lip}}
\DeclareMathOperator*{\argmax}{arg\,max}
\DeclareMathOperator*{\argmin}{arg\,min}
\endlocaldefs

\let\NBOLongtable\longtable
\def\longtable{\Needspace{7\baselineskip}\NBOLongtable}

\begin{document}

\begin{frontmatter}

\title{Response to the Referee: Neural Bellman Operators}

\runtitle{Response to the Referee}

\begin{aug}

\author[id=au1,addressref={add1}]{\fnms{Qian}~\snm{QI}}

\address[id=add1]{Peking University}

\end{aug}

\end{frontmatter}

\noindent Revision R21, 5 October 2026. This response addresses the advisory report on the R18 manuscript. The report was commissioned by the repository owner, not by the Econometric Society.

\par We thank the referee for distinguishing valid local continuation evidence from broader economic claims. We strengthen the theorem-to-policy link, integrate the completed future-change experiment, and report certified economic consequences. The original title, economic applications, and unfavorable evidence are retained. The experiments below are inherited frozen executions; their replay and new reporting are not counted as new independent observations.

\subsection*{B1. The positive result and the class of changing economic futures}

We retain the original controlled-economy problem and title. The new section on policy accuracy proves a monotone-operator composition result for certificates evaluated against the implemented policy's own future, together with a state-dependent additive envelope and an explicit class/value-transfer bridge. Its uniform-domain and information assumptions are stated, not inferred from a finite sample. This extends the analytical connection to the original recursive and strategic applications without asserting that the finite capital experiment numerically verifies those global premises. The completed, previously frozen future-change study now changes subsequent policy, technology, and payoff weights: it is not another current-task menu with an identical future. All 168 services and 840 outcomes are incorporated.


\par\smallskip\noindent\emph{Location.} Sections on policy accuracy and changing futures; complete proofs and records in the supplement.

\subsection*{B2. Symmetric, decision-relevant prediction diagnostics}

The R19 common-action and own-action design is retained. The integrated R20 record evaluates every predictor at exogenous lower, midpoint, and upper actions and separately at its own selected actions against exact finite-law continuation means. Absolute risk, centered risk, own-action risk, and actual three-action decision loss remain separate, with caches rebuilt over three newly declared streams. A particularly informative comparison is the large absolute risk after lowering future withdrawal despite very small centered risk and decision loss. These diagnostics are not substituted for full-interval policy accuracy: the actual stopping test uses the true objective and its curvature. The earlier conditional 64-path Raw comparisons remain unchanged.


\par\smallskip\noindent\emph{Location.} Changing-futures mechanism subsection; Tables of common-action and own-action risks in the supplement.

\subsection*{B3. Prospective stopping rather than retrospective selected work}

The earlier exact catalogue closure remains a retrospective use of simultaneous intervals. It has not been relabeled as an executed stopping procedure. The R19 and R20 designs instead have frozen stage sequences and first-success checks, with every failed attempt charged. The R20 service either reuses, refreshes under its registered rule, refits under its own registered rule, or records an exhausted budget. No procedure is selected after looking at outcomes and then reported as though only its cost had been paid. Deterministic finite-law interval checks are valid under stopping without statistical error spending. The original complete comparative bill, 1,897.326 seconds, is reported alongside the deployed-service components.


\par\smallskip\noindent\emph{Location.} Work-to-tolerance section; changing-futures work table and chronology; immutable original records.

\subsection*{B4. The measured amortization frontier}

The six-volume R19 experiment, at 1, 4, 16, 64, 256, and 1,024 queries, is preserved. The changing-futures experiment adds actual five-future services at 1, 16, 256, and 1,024 queries per future. The paper reports measured initialization, fitting, querying, verification, failed checks, and record costs without fitting a line through one workload. The new complete table shows that adaptive NBO is more expensive than fresh-refit NBO in every displayed dimension-volume cell; we explain the anchor and failed-check costs and retain them. The sufficient amortization inequality and the new across-date allocation result characterize explicit cost accounts, not an unsupported universal break-even point.


\par\smallskip\noindent\emph{Location.} Existing six-volume tables; new eight-cell cost table; work-allocation proposition.

\subsection*{B5. Contribution-isolating baselines and numerical antecedents}

The quadratic and radial-basis postdecision regressions, shared task-conditioned actor, cached SAA, and full-support classical method from R19 remain. R20 applies the same fitting and selection implementations to genuinely changed futures and includes both conventional adaptive surrogates. Their successes and competitive clocks are printed, not relegated to an omitted robustness exercise. The literature discussion explicitly distinguishes approximate policy iteration, continuation regression, actor/gradient methods, simulation optimization, and standard certificate constructions. The new composition proof is identified as a dynamic-programming stability argument; its role is the economic verification chain, not a new Bellman principle. We add the primary finite-time fitted-value reference and explain why sampling norms do not alone imply the uniform policy premise. No claim of an exhaustive comparison against every possible surrogate is made.


\par\smallskip\noindent\emph{Location.} Literature discussion, policy-composition section, and both prospective comparisons.

\subsection*{B6. An economic consequence rather than only a ranking}

The new economic contrast uses NBO's own saved actions and derivative and curvature certificates to enclose the true mean optimal current-withdrawal response to stronger future production and greater future valuation. All four dimension-by-change upper endpoints are negative. For example, the fifty-dimensional production response is enclosed by [-0.037119,-0.026553] in withdrawal units. The table reports outward-rounded bounds rather than signs from noisy point estimates. Enumeration independently confirms and sharpens the target; it is not passed off as the neural certificate. We retain the task-specific external log-utility compensation interpretation and explicitly avoid exponentiating a mean loss bound into a bound on mean compensation. The exercise is a declared capital counterfactual, not an empirical tax estimate.


\par\smallskip\noindent\emph{Location.} New NBO economic-contrast table; proof of the optimal-action radius; enumeration benchmark.

\subsection*{B7. Finite decisions, dynamic policy loss, and continuous-time HJB}

The policy-composition theorem now states the exact missing link: uniform full-action advantages against the implemented policy's own continuation propagate with the evaluation operators' Lipschitz weights. A two-date counterexample shows why an observed catalogue or the implemented policy's occupancy distribution cannot replace state coverage. The continuous-economy bridge adds class approximation, value transfer, and implementation allowances explicitly. These are conditions for a general guarantee, not numerical allowances supplied by the finite mean-loss experiment. The original diffusion verification, strengthened-HJB results, unresolved comparisons, and full application models remain in their own accounts.


\par\smallskip\noindent\emph{Location.} Policy-composition theorem, continuous-economy bridge, supplement counterexample, retained continuous-time evidence.

\subsection*{B8. A coherent main argument without deleting the paper}

We keep Neural Bellman Operators, the general economic problem, and the complete application companion. Six entire historical theory/design/risk/catalogue sections are relocated intact to the technical supplement, while the main article follows the decision target, certificates, original controlled model, and executed economic comparisons. No historical component is deleted; original roots and changed introduction/conclusion are archived, and an automated source-and-label audit checks retention. The new future-change evidence is integrated into the paper itself rather than left as an unconnected code branch. This addresses readability by organization, not by abandoning the original subject.


\par\smallskip\noindent\emph{Location.} New root article and supplement; retained-section guide; inheritance and preservation audit.

\subsection*{M1. Center the economic decision without changing the subject}

The decision problem, economic target, and complete-work meaning appear before the general controlled model. The title and broad economic applications are retained. The paper presents the common evaluation-decision-payoff argument rather than requiring a reader to navigate branch history. Revision identifiers are kept in the reproducibility record, not used as substitutes for economic explanation.


\par\smallskip\noindent\emph{Location.} Introduction, target hierarchy, and root ordering.

\subsection*{M2. Freeze, execute, and charge the stopping procedure}

The R19 prospective design and R20 future-change source freeze are preserved with their actual execution records. Candidate stages, tolerance, checks, and failure treatment precede the execution. R21's new tables and optimal-response bands are expressly labeled deterministic post-execution analyses of those saved observations, not a backdated new registration. There are no invented new fits or replacement timings.


\par\smallskip\noindent\emph{Location.} Chronology, source manifest, full service record.

\subsection*{M3. Report query-volume and component work}

Both measured volume grids and the complete accounting boundaries are reported. A service may be compared only with accuracy-qualified alternatives at the same task volume. Cheap unresolved reuse is visible but is not an accuracy-qualified winner; shared startup and post-stop diagnostics remain on the complete comparison bill.


\par\smallskip\noindent\emph{Location.} Six-volume inherited study and new eight-cell future-service table.

\subsection*{M4. Report common-action, own-action, and decision diagnostics}

The complete predictor diagnostics use exact finite-law means at common exogenous actions and each predictor's selected actions. Cache randomness is regenerated across three fixed streams. The target remains conditional on the declared finite economic laws; these three streams are not advertised as establishing a population guarantee over all future training runs.


\par\smallskip\noindent\emph{Location.} Mechanism discussion and full risk table.

\subsection*{M5. Use strong conventional continuation baselines}

The same-data quadratic and radial-basis regressions, shared direct actor, cached SAA, and structure-exploiting enumerated calculation remain in the comparison. Their successful and cheaper outcomes are retained. The changed-future experiment also compares adaptive conventional surrogates with adaptive and fresh neural continuation fitting.


\par\smallskip\noindent\emph{Location.} Method definitions and complete results in both study families.

\subsection*{M6. Make the economic counterfactual explicit}

The temporary-charge exercise is joined by changes in future withdrawal, production coupling, and valuation. NBO action certificates establish signed current-response bands. The model units and task-specific welfare interpretation are explicit; no external calibration or observed reform is invented.


\par\smallskip\noindent\emph{Location.} Economic design and certified response subsection.

\subsection*{M7. Identify standard mathematics correctly}

Bias-variance identities, common-noise cancellation, concentration bounds, shortest-path closure, Bellman stability, and the convex budget allocation are not advertised as new underlying mathematical principles. The paper's contribution is the particular economic certificate chain, its reusable continuation implementation, and the complete measured evidence under specified information and work accounts.


\par\smallskip\noindent\emph{Location.} Literature and policy-accuracy discussion; proofs retain explicit antecedents.

\subsection*{M8. Maintain the hierarchy of accuracy targets}

The main hierarchy table is retained. Every new table distinguishes prediction risk, finite-menu decision loss, mean full-scalar-interval regret, optimal-action response, and work. The policy theorem and continuous-economy bridge require their separate uniform and transfer premises. No finite catalogue result is used as a numerical full-HJB certificate.


\par\smallskip\noindent\emph{Location.} Target hierarchy, new table captions, composition assumptions.

\subsection*{M9. Retain all adverse and unresolved findings}

All five specific R18 cautions remain: cheaper Raw/vector final clocks, SAA advantages in the Long cells, one least-cost-certified NBO cell, no favorable 64-path risk comparison, and cheaper unresolved candidates in the decisive cell. R21 additionally prints all 48 uncertified unchanged-future outcomes and the adaptive-versus-refit cost disadvantage. These facts define useful operating regions rather than being removed to improve a ranking.


\par\smallskip\noindent\emph{Location.} Introduction, retained historical accounts, full future accuracy and cost tables.

\subsection*{M10. Improve structure while preserving complete material}

The main reading sequence is reorganized and complete sections are relocated, not shortened by deleting results or changing the topic. All economic applications and proof blocks remain available, and every inherited label is checked. The article, technical supplement, applications, and this response are compiled in the native econsocart class with synchronized references.


\par\smallskip\noindent\emph{Location.} Four native-class publication files, source-preservation map, and compilation audit.

\subsection*{Verification and scope}

The delivery contains new unit tests, the inherited test suite, a deterministic replay of every saved future-change certificate, a source/label preservation audit, and four native-class PDF builds. Exact results and PDF hashes are recorded in the accompanying audit files. The new frozen-evidence integration does not constitute a new prospective experiment. Uniform state coverage for the general policy theorem and continuous-economy transfer allowances must be supplied when using that theorem in a particular application; the finite service study does not supply them by implication. No assertion of journal acceptance or universal neural superiority is made.

\end{document}
